\documentclass[11pt,a4paper]{article}

\usepackage[a4paper,margin=1in]{geometry}
\usepackage{fontspec}
\IfFontExistsTF{Times New Roman}
  {\setmainfont{Times New Roman}}
  {\IfFontExistsTF{TeX Gyre Termes}{\setmainfont{TeX Gyre Termes}}{\setmainfont{Times}}}
\usepackage{newtxmath}
\usepackage{amsmath}
\usepackage{booktabs}
\usepackage{longtable}
\usepackage{array}
\usepackage{enumitem}
\usepackage{xcolor}
\usepackage{listings}
\usepackage[round,authoryear]{natbib}
\usepackage[hidelinks]{hyperref}
\usepackage{xurl}
\usepackage{fancyhdr}
\usepackage{microtype}

\setlength{\parindent}{0pt}
\setlength{\parskip}{0.65em}
\setlist[itemize]{leftmargin=1.4em,itemsep=0.2em,topsep=0.2em}
\setlist[enumerate]{leftmargin=1.6em,itemsep=0.2em,topsep=0.2em}

\pagestyle{fancy}
\fancyhf{}
\lhead{\small Ticket 03 Study Report}
\rhead{\small European CN Validation}
\cfoot{\thepage}
\renewcommand{\headrulewidth}{0.3pt}

\lstset{
  basicstyle=\ttfamily\small,
  breaklines=true,
  frame=single,
  columns=fullflexible,
  keepspaces=true,
  rulecolor=\color{black!30},
  xleftmargin=0.5em,
  xrightmargin=0.5em
}

\newcommand{\code}[1]{\texttt{#1}}
\newcommand{\file}[1]{\path{#1}}
\newcolumntype{L}[1]{>{\raggedright\arraybackslash}p{#1}}
\newcolumntype{P}[1]{>{\raggedright\arraybackslash\ttfamily\footnotesize}p{#1}}

\title{\textbf{Ticket 03 Study Report: European Crank-Nicolson Validation}\\
\large SURF2026 American Option Risk Surfaces}
\author{Codex-assisted implementation report}
\date{June 15, 2026}

\begin{document}
\maketitle

\begin{abstract}
Ticket 03 implements a European-only Crank-Nicolson validation mode for the one-dimensional
Black-Scholes PDE with continuous dividends. It uses the Ticket 01 payoff and closed-form
Black-Scholes utilities, plus the Ticket 02 grid and spatial-operator helpers. The goal is to
check the finite-difference time-stepping machinery against known European option prices before
adding any American payoff obstacle, projected SOR iteration, free-boundary extraction, Greeks,
stress maps, datasets, or neural-network work. This report is written for beginner researchers
with accounting, data, machine-learning, and Python experience who are still learning PDE-based
option pricing.
\end{abstract}

\tableofcontents
\newpage

\section{Role of Ticket 03 in the Whole Solver Project}

The project is building toward validated American option risk surfaces. The final solver will
eventually need to handle early exercise, an obstacle condition, a continuation region, an
exercise region, and a free boundary. Before those harder American-option features are added,
the project needs evidence that the simpler European PDE machinery is correct.

Ticket 03 is that checkpoint. It implements Crank-Nicolson time stepping for European calls and
puts only. European options are valuable for validation because, under the Black-Scholes model
with continuous dividends, we already have closed-form prices from Ticket 01. A finite-difference
solver can therefore be compared directly against analytic benchmark values.

The logic of the ticket order is:
\begin{itemize}
  \item Ticket 01 builds payoff and closed-form Black-Scholes utilities.
  \item Ticket 02 builds the uniform grids, interior node convention, operator coefficients,
        and boundary helpers.
  \item Ticket 03 combines those pieces into European Crank-Nicolson time stepping and checks the
        numerical prices against closed-form European prices.
  \item Ticket 04 can later add American obstacle and PSOR logic only after this European
        continuation mode is credible.
\end{itemize}

This ticket is deliberately not a full solver validation report. It validates one layer: European
continuation time stepping on top of the existing grid and operator setup. That narrowness is a
strength, because it makes numerical mistakes easier to diagnose before the American problem adds
more moving parts.

\section{Theory and Methodology Connection}

The source basis for this ticket is internal and project-specific:
\begin{itemize}
  \item \file{reports/01_literature/literature_map.md};
  \item \file{reports/02_math/formulation_note.md};
  \item \file{reports/03_solver/solver_validation_plan.md};
  \item \file{reports/03_solver/solver_implementation_tickets.md};
  \item the Ticket 01 and Ticket 02 study reports and code.
\end{itemize}

The literature map motivates a classical finite-difference benchmark before neural surrogate
modelling. The formulation note defines the dividend-adjusted Black-Scholes operator. The solver
validation plan identifies European closed-form comparison as the first numerical validation case.
This reference pass connects that implementation to the original Crank-Nicolson time-stepping
method \citep{crank_nicolson_1947}, the Black-Scholes/Merton option-pricing framework
\citep{black_scholes_1973,merton_1973}, and standard finite-difference option-pricing background
\citep{wilmott_howison_dewynne_1995}.

The continuous Black-Scholes spatial operator is
\begin{equation}
  L U =
  \frac{1}{2}\sigma^2 S^2 U_{SS}
  + (r-q)S U_S
  - rU.
  \label{eq:operator}
\end{equation}

Here \(S\) is spot, \(\sigma\) is volatility, \(r\) is the risk-free rate, and \(q\) is the
continuous dividend yield. In time-to-maturity notation, the European continuation PDE is
\begin{equation}
  U_{\tau} = L U,
  \qquad U(S,0)=\Phi(S),
  \label{eq:continuation-pde}
\end{equation}
where \(\Phi\) is the call or put payoff. The solver starts from the known payoff at
\(\tau=0\) and marches forward to \(\tau=T\).

\section{Why European PDE Validation Comes Before American CN/PSOR}

European options do not have an early-exercise right. That means there is no payoff obstacle and
no free boundary. The numerical problem is an ordinary linear continuation PDE. This makes
European validation the cleanest way to test:
\begin{itemize}
  \item whether time-to-maturity is marched in the correct direction;
  \item whether the finite-difference operator sign convention matches the formulation note;
  \item whether continuous dividends enter through \(r-q\) and the closed-form benchmark;
  \item whether boundary values are applied consistently at old and new time levels;
  \item whether the tridiagonal linear solve is working.
\end{itemize}

American CN/PSOR will later add a projection step that keeps the option value above payoff.
If the European continuation step is wrong, adding projection would make debugging harder rather
than easier. A projected method might hide an error in the continuation equation by forcing values
back above payoff. Ticket 03 therefore validates the continuation engine before any projection is
introduced.

For beginner researchers, the key idea is this: European validation checks the ``PDE engine.''
American validation will later check the ``PDE engine plus exercise constraint.'' The project
should not mix those two questions too early.

\section{Crank-Nicolson in Plain Language}

A finite-difference method replaces continuous spot and time with grid values. Ticket 02 already
created the spot grid \(S_i\), the time-to-maturity grid \(\tau_n\), and the discrete spatial
operator \(L_h\). Ticket 03 adds the time-stepping rule.

The Crank-Nicolson method averages the PDE operator between the old time level and the new time
level:
\begin{equation}
  \frac{U^{n+1} - U^n}{\Delta \tau}
  =
  \frac{1}{2}L_h U^{n+1}
  +
  \frac{1}{2}L_h U^n.
  \label{eq:cn-average}
\end{equation}

In plain language, instead of using only the old slope or only the new slope, Crank-Nicolson uses
an average of both. This usually gives better accuracy than a purely explicit first-order method,
while still requiring a linear system solve at each time step.

Rearranging Equation~\ref{eq:cn-average} gives
\begin{equation}
  \left(I - \frac{\Delta \tau}{2}L_h\right)U^{n+1}
  =
  \left(I + \frac{\Delta \tau}{2}L_h\right)U^n.
  \label{eq:cn-matrix}
\end{equation}

The implementation solves only the interior spot nodes. Boundary nodes are imposed separately
using the European boundary helpers from Ticket 02.

\section{Matrix Form and Boundary Adjustments}

Ticket 02 represents the spatial operator at interior node \(i\) as
\begin{equation}
  (L_h U)_i = a_i U_{i-1} + b_i U_i + c_i U_{i+1},
  \label{eq:lh-row}
\end{equation}
where \(a_i\), \(b_i\), and \(c_i\) are the lower, diagonal, and upper coefficients.

For Ticket 03, define
\begin{equation}
  A = I - \frac{\Delta \tau}{2}L_h,
  \qquad
  b^n = \left(I + \frac{\Delta \tau}{2}L_h\right)U^n.
\end{equation}

On the old time level \(n\), the boundary values are already part of the full vector \(U^n\), so
they naturally enter the expression \(L_h U^n\). On the new time level \(n+1\), the unknown vector
contains only interior values, so boundary contributions must be moved to the right-hand side.
For the first and last interior rows, the adjustment is
\begin{align}
  b^n_1 &\leftarrow b^n_1
    + \frac{\Delta \tau}{2} a_1 B^{n+1}_0, \\
  b^n_{M-1} &\leftarrow b^n_{M-1}
    + \frac{\Delta \tau}{2} c_{M-1} B^{n+1}_M,
\end{align}
where \(B^{n+1}_0\) and \(B^{n+1}_M\) are the imposed lower and upper boundary values at the new
time level. This sign is easy to get wrong. The plus sign appears because the left-hand matrix has
boundary coefficients \(-\frac{\Delta \tau}{2}a_1\) and
\(-\frac{\Delta \tau}{2}c_{M-1}\), and moving those known terms to the right-hand side changes the
sign.

The linear system is tridiagonal, so Ticket 03 uses a small Thomas algorithm helper instead of
SciPy. This keeps the implementation dependency-light and consistent with Tickets 01 and 02.

\section{Implementation Design Decisions}

\subsection{Use Existing Ticket 01 and Ticket 02 Utilities}

The solver reuses Ticket 01 payoff and closed-form pricing functions. It also reuses Ticket 02
spot grids, time grids, spatial operator coefficients, operator application, and European boundary
helpers. This keeps Ticket 03 focused on time stepping rather than recreating earlier layers.

\subsection{European-Only Public API}

The public solver function accepts only \code{option\_type="call"} or \code{"put"}. It rejects
other option types with \code{ValueError}. No American projection, obstacle enforcement, free-
boundary extraction, Greeks, stress maps, datasets, or neural-network files are introduced.

\subsection{NumPy and Standard Library Only}

SciPy is avoided for this ticket. The tridiagonal linear system is solved with a local Thomas
algorithm helper. The validation CSV is written with the standard-library \code{csv} module rather
than \code{pandas}. This matches the dependency caution recorded in earlier tickets.

\subsection{Target Region Error Metrics}

The validation plan emphasizes the region \(S/K \in [0.4,1.8]\), because that is the likely
interpretation region for later risk surfaces. Ticket 03 reports maximum absolute error, RMSE,
and the spot location of the maximum error in this target region. Boundary-region errors are not
hidden, but they are not the headline metric for this first validation table.

\subsection{No Optional Figure}

No figure is generated in Ticket 03. The required numerical validation table is enough for this
checkpoint, and avoiding \code{matplotlib} keeps the artifact generation simple and reproducible
across the current Python runtimes.

\section{Code-Level Function Explanations}

\begin{longtable}{P{0.37\textwidth}L{0.55\textwidth}}
\toprule
Function or class & Purpose and important behavior \\
\midrule
\endhead
\code{ValidationMetrics} &
Frozen dataclass holding max absolute error, RMSE, max-error spot, and the target moneyness bounds. \\
\code{EuropeanCNResult} &
Frozen dataclass bundling option parameters, grids, numerical values, closed-form values, errors, and metrics. \\
\code{solve\_tridiagonal} &
Thomas-algorithm helper for tridiagonal systems. It validates array dimensions and rejects zero pivots. \\
\code{target\_region\_error\_metrics} &
Computes max absolute error, RMSE, and max-error spot over \(S/K\in[0.4,1.8]\) by default. \\
\code{european\_crank\_nicolson\_price} &
Builds grids, initializes payoff at \(\tau=0\), applies European CN time stepping, compares against closed-form Black-Scholes prices, and returns a result dataclass. \\
\file{experiments/01_european_cn_validation.py} &
Runs default put/call medium and fine grid cases, then writes the validation CSV. \\
\bottomrule
\end{longtable}

Every new Python source file created in this ticket states in its module docstring that it was
created for Ticket 03.

\section{Testing Strategy}

\begin{longtable}{L{0.25\textwidth}L{0.31\textwidth}L{0.32\textwidth}}
\toprule
Test category & What it proves & What it does not prove \\
\midrule
\endhead
Tridiagonal solver test &
The local Thomas algorithm solves a known small system. &
All possible matrix conditioning cases. \\
\(T=0\) payoff test &
The solver respects the maturity condition and returns payoff without time stepping. &
Positive-maturity PDE accuracy. \\
European put CN comparison &
Put finite-difference values are close to Ticket 01 closed-form values in a representative case. &
American put obstacle correctness. \\
European call CN comparison &
Call finite-difference values are close to Ticket 01 closed-form values in a representative dividend case. &
No-dividend American call validation. \\
Error metric test &
Max absolute error, RMSE, and max-error location are computed correctly on a known vector. &
Solver accuracy itself. \\
Invalid input tests &
Bad option type, strike, maturity, volatility, domain, and grid sizes raise clear errors. &
A full production validation policy. \\
No PSOR/API test &
The Ticket 03 public module surface does not introduce PSOR or American projection names. &
Future Ticket 04 behavior. \\
\bottomrule
\end{longtable}

The tests are intentionally focused. They are designed to catch sign errors, boundary mistakes,
bad linear-solver behavior, and accidental scope creep. They do not claim formal convergence or
complete production robustness.

\section{Validation Result Table and Interpretation}

The validation script wrote:
\begin{lstlisting}
results/01_solver_validation/tables/ticket_03_european_cn_validation.csv
\end{lstlisting}

The table below reports errors over the target region \(S/K\in[0.4,1.8]\), with \(K=1\),
\(T=1\), \(r=0.05\), \(q=0.02\), \(\sigma=0.20\), and \(S_{\max}=4\).

\begin{longtable}{L{0.13\textwidth}r r r r}
\toprule
Option & \(M\) & \(N\) & Max absolute error & RMSE \\
\midrule
\endhead
Put & 120 & 120 & \(2.68274272102\times10^{-4}\) & \(1.15854293529\times10^{-4}\) \\
Put & 180 & 180 & \(1.19448913490\times10^{-4}\) & \(5.16267098347\times10^{-5}\) \\
Call & 120 & 120 & \(2.68273627869\times10^{-4}\) & \(1.15853941701\times10^{-4}\) \\
Call & 180 & 180 & \(1.19448627389\times10^{-4}\) & \(5.16265526360\times10^{-5}\) \\
\bottomrule
\end{longtable}

The largest errors for these cases occurred near \(S=0.9667\) on the \(120\times120\) grid and
near \(S=0.9778\) on the \(180\times180\) grid. That location is close to the payoff kink around
\(S=K=1\), which is expected for a plain Crank-Nicolson start from a nonsmooth payoff.

The medium-to-fine comparison improves for both calls and puts. The maximum absolute error drops
from approximately \(2.68\times10^{-4}\) to \(1.19\times10^{-4}\), and RMSE drops from about
\(1.16\times10^{-4}\) to \(5.16\times10^{-5}\). This is useful validation evidence for the
European continuation step.

This does not validate the future American solver. There is no obstacle projection, no PSOR
iteration, no complementarity residual, no boundary extraction, and no Greek calculation in this
ticket.

\section{Limitations}

Ticket 03 has deliberate limits:
\begin{itemize}
  \item no PSOR;
  \item no American option solving;
  \item no American payoff obstacle or projection;
  \item no free-boundary extraction;
  \item no Greeks;
  \item no stress maps;
  \item no datasets;
  \item no neural-network files;
  \item no optional figure;
  \item no Rannacher smoothing;
  \item no formal convergence proof.
\end{itemize}

The current validation cases are representative, not exhaustive. They use one parameter set with
continuous dividends and two grid sizes. Later validation should include more maturities,
volatilities, dividend yields, grid refinements, and domain sensitivity checks before treating the
solver as a trusted label generator.

\section{Lessons Learned / Study Notes}

The first lesson is that time direction matters. In time-to-maturity form, the solver starts from
the payoff at \(\tau=0\) and marches forward to \(\tau=T\). This can feel backwards if one is used
to calendar time, but it is natural for option-pricing finite differences.

The second lesson is that Crank-Nicolson is a bridge between PDEs and linear algebra. The PDE
\(U_{\tau}=LU\) becomes a tridiagonal system at each time step. Understanding the matrix form
helps explain why a linear solver is needed even before American projection appears.

The third lesson is that boundary values enter twice: old boundary values are part of \(U^n\), and
new boundary values must be adjusted into the right-hand side. Many beginner solver bugs come from
forgetting one of these boundary contributions or using the wrong sign.

The fourth lesson is that closed-form European prices are a powerful benchmark. They allow us to
test the numerical engine before early exercise makes the problem harder.

The fifth lesson is to avoid overclaiming. Passing Ticket 03 means the European CN validation mode
works for the tested cases. It does not mean that American exercise, PSOR convergence,
free-boundary extraction, or Greeks are correct.

\section{Link to Ticket 04}

Ticket 04 should add the American obstacle and PSOR/LCP core. It can reuse the Ticket 03
continuation time step structure, but it must change the solve at each time step: instead of
solving one unconstrained European linear system, it must enforce the American condition
\begin{equation}
  U^{n+1} \geq \Phi.
\end{equation}

The correct next step is therefore not stress-map generation or neural modelling. The correct next
step is to implement and test the projected linear complementarity solve for American options,
with convergence metadata and obstacle diagnostics.

\section{Reproducibility Record}

\subsection{Files Created}

\begin{longtable}{P{0.52\textwidth}L{0.40\textwidth}}
\toprule
Path & Role in Ticket 03 \\
\midrule
\endhead
\file{src/american_risk_surfaces/solvers/cn.py} &
European CN validation solver, tridiagonal helper, result dataclasses, and target-region metrics. \\
\file{tests/test_cn_european.py} &
Focused \code{unittest} coverage for Ticket 03 behavior. \\
\file{experiments/01_european_cn_validation.py} &
Reproducible CSV validation runner. \\
\file{results/01_solver_validation/tables/ticket_03_european_cn_validation.csv} &
Generated numerical validation table. \\
\file{reports/03_solver/tickets/ticket_03_european_cn_validation.tex} &
Formal LaTeX study-report source. \\
\file{reports/03_solver/tickets/ticket_03_european_cn_validation.pdf} &
Formal PDF compiled from the LaTeX source. \\
\bottomrule
\end{longtable}

\subsection{Commands}

Test command:
\begin{lstlisting}[language=bash]
PYTHONPYCACHEPREFIX=/private/tmp/codex_pycache PYTHONPATH=src python3 -m unittest discover -s tests
\end{lstlisting}

Final observed test result:
\begin{lstlisting}
Ran 32 tests in 0.024s

OK
\end{lstlisting}

Validation CSV command:
\begin{lstlisting}[language=bash]
PYTHONPYCACHEPREFIX=/private/tmp/codex_pycache PYTHONPATH=src python3 experiments/01_european_cn_validation.py
\end{lstlisting}

Python compile command:
\begin{lstlisting}[language=bash]
PYTHONPYCACHEPREFIX=/private/tmp/codex_pycache PYTHONPATH=src python3 -m compileall -q src tests experiments
\end{lstlisting}

LaTeX compile command:
\begin{lstlisting}[language=bash]
HOME=/private/tmp FC_CACHEDIR=/private/tmp/fontconfig-cache latexmk -xelatex -interaction=nonstopmode -halt-on-error -outdir=/private/tmp/ticket03_latex reports/03_solver/tickets/ticket_03_european_cn_validation.tex
\end{lstlisting}

PDF verification uses \code{pdfinfo} and page rendering with \code{pdftoppm}. Raw Git status is
intentionally not included in this formal PDF report.

\bibliographystyle{plainnat}
\bibliography{reports/03_solver/references}

\end{document}
